\ifdefined\gkver\else\def\gkver{student}\fi
\documentclass[\gkver]{gaokaozhenti}
\begin{document}

\chapter{1991年全国}

\begin{enumerate}
\item
%% number: 1
%% paperName: 1991年全国高考第 1 题 2 分
%% typeId: 1
%% type: 单选题
%% chapter: 直线运动
%% point: 竖直上下抛运动
%% method: 无
%% score: 2
%% degree: 850
%% duplicateId: 0
%% body:
以初速 $v_{0}$ 竖直上抛一小球。若不计空气阻力，在上升过程中，从抛出到小球动能减少一半所经过的时间是\xzanswer{D}
\fourchoices[answer=D]
{$\frac{v_{0}}{g}$}
{$\frac{v_{0}}{2g}$}
{$\frac{\sqrt{2}v_{0}}{2g}$}
{$\frac{(2-\sqrt{2})v_{0}}{2g}$}

%% answer: D
%% memo:
\memoanswer{
小球做竖直上抛运动，由题意可得 $\frac{1}{2}mv_{t}^{2}=\frac{1}{2}\cdot\frac{1}{2}mv_{0}^{2}$，解得 $v_{t}=\frac{\sqrt{2}}{2}v_{0}$，由运动学公式得 $v_{t}=v_{0}-gt$，解得 $t=\frac{v_{0}}{g}\left(1-\frac{\sqrt{2}}{2}\right)$，故 D 正确。
}

\item
%% number: 2
%% paperName: 1991年全国高考第 2 题 2 分
%% typeId: 1
%% type: 单选题
%% chapter: 电场
%% point: 带电粒子的直线运动
%% method: 无
%% score: 2
%% degree: 850
%% duplicateId: 0
%% body:
下列粒子从初速为零的状态经过加速电压为 $U$ 的电场之后，哪种粒子的速度最大？\xzanswer{A}
\fourchoices[answer=A]
{质子}
{氘核}
{$\alpha$ 粒子}
{钠离子 \ce{Na+}}

%% answer: A
%% memo:
\memoanswer{
由动能定理得 $qU=\frac{1}{2}mv^{2}-0$，解得 $v=\sqrt{\frac{2qU}{m}}$，可知 $U$ 一定时，$v$ 与 $\sqrt{\frac{q}{m}}$ 成正比，$\frac{q}{m}$ 越大，$v$ 也越大，4 种粒子的比荷分别是：①质子：1，②氘核：$\frac{1}{2}$，③$\alpha$ 粒子：$\frac{1}{2}$，④Na 离子：$\frac{1}{23}$，故 A 正确。
}

\item
%% number: 3
%% paperName: 1991年全国高考第 3 题 2 分
%% typeId: 1
%% type: 单选题
%% chapter: 磁场
%% point: 磁力矩
%% method: 无
%% score: 2
%% degree: 750
%% duplicateId: 0
%% body:
如图所示，一位于 $xy$ 平面内的矩形通电线圈只能绕 $Ox$ 轴转动，线圈的四个边分别与 $x$、$y$ 轴平行。线圈中电流方向如图。当空间加上如下所述的哪种磁场时，线圈会转动起来？\xzanswer{B}
\onepicture[w=4cm]{\includesvg{figs/03}}
\fourchoices[answer=B]
{方向沿 $x$ 轴的恒定磁场}
{方向沿 $y$ 轴的恒定磁场}
{方向沿 $z$ 轴的恒定磁场}
{方向沿 $z$ 轴的变化磁场}

%% answer: B
%% memo:
\memoanswer{
线圈只能绕 $Ox$ 轴转动，表明图示位置时 $ab$ 边与 $dc$ 边分别受到平行于 $z$ 轴方向相反的两个力作用，由左手定则知磁场方向与 $y$ 轴同向或反向，故 B 正确。
}

\item
%% number: 4
%% paperName: 1991年全国高考第 4 题 1 分
%% typeId: 1
%% type: 单选题
%% chapter: 功和能
%% point: 功率
%% method: 无
%% score: 1
%% degree: 750
%% duplicateId: 0
%% body:
一质量为 $m$ 的木块静止在光滑的水平面上。从 $t = 0$ 开始，将一个大小为 $F$ 的水平恒力作用在该木块上。在 $t = t_{1}$ 时刻力 $F$ 的功率是\xzanswer{C}
\fourchoices[answer=C]
{$\frac{F^{2}}{2m}t_{1}$}
{$\frac{F^{2}}{2m}t_{1}^{2}$}
{$\frac{F^{2}}{m}t_{1}$}
{$\frac{F^{2}}{m}t_{1}^{2}$}

%% answer: C
%% memo:
\memoanswer{
由功率定义知 $p=Fv=Fat_{1}=F\frac{F}{m}t_{1}=\frac{F^{2}}{m}t_{1}$，故 C 正确。
}

\item
%% number: 5
%% paperName: 1991年全国高考第 5 题 2 分
%% typeId: 1
%% type: 单选题
%% chapter: 曲线运动
%% point: 平抛运动
%% method: 无
%% score: 2
%% degree: 650
%% duplicateId: 0
%% body:
如图所示，以 9.8 米/秒的水平初速度 $v_{0}$ 抛出的物体，飞行一段时间后，垂直地撞在倾角 $\theta$ 为 $30^{\circ}$ 的斜面上。可知物体完成这段飞行的时间是\xzanswer{C}
\onepicture[w=5cm]{\includesvg{figs/05}}
\fourchoices[answer=C]
{$\frac{\sqrt{3}}{3}$ s}
{$\frac{2\sqrt{3}}{3}$ s}
{$\sqrt{3}$ s}
{2 s}

%% answer: C
%% memo:
\memoanswer{
将垂直于斜面的速度 $v$ 分解，得 $v_y = \frac{v_0}{\tan 30^{\circ}} = \sqrt{3}v_0$，由 $v_y = gt$ 得 $t = \sqrt{3}\Us$，故 C 正确。
}

\item
%% number: 6
%% paperName: 1991年全国高考第 6 题 2 分
%% typeId: 1
%% type: 单选题
%% chapter: 运动和力
%% point: 牛顿第二定律
%% method: 无
%% score: 2
%% degree: 750
%% duplicateId: 0
%% body:
有两个物体 a 和 b，其质量分别为 $m_{a}$ 和 $m_{b}$，且 $m_{a} > m_{b}$。它们的初动能相同。若 a 和 b 分别受到不变的阻力 $F_{a}$ 和 $F_{b}$ 的作用，经过相同的时间停下来，它们的位移分别为 $s_{a}$ 和 $s_{b}$，则\xzanswer{A}
\fourchoices[answer=A]
{$F_{a} > F_{b}$ 且 $s_{a} < s_{b}$}
{$F_{a} > F_{b}$ 且 $s_{a} > s_{b}$}
{$F_{a} < F_{b}$ 且 $s_{a} > s_{b}$}
{$F_{a} < F_{b}$ 且 $s_{a} < s_{b}$}

%% answer: A
%% memo:
\memoanswer{
由题意知，$\frac{1}{2}m_a v_a^2 = \frac{1}{2}m_b v_b^2$，又 $m_a > m_b$，所以 $v_a < v_b$，由位移公式知 $s = \bar{v} \cdot t = \frac{v_0 - 0}{2} \cdot t$，因 $t$ 相同，所以 $s \propto v_0$，所以 $s_a < s_b$①，由于两物体的动能变化相同，由动能定理可知外力做功相等，即 $F_a s_a = F_b s_b$②，将①代入②得 $F_a > F_b$，故 A 正确。
}

\item
%% number: 7
%% paperName: 1991年全国高考第 7 题 2 分
%% typeId: 1
%% type: 单选题
%% chapter: 力物体的平衡
%% point: 力矩平衡
%% method: 无
%% score: 2
%% degree: 450
%% duplicateId: 0
%% body:
图中 A、B 是两块相同的均匀长方形砖块，长为 $l$，叠放在一起，A 砖相对于 B 砖右端伸出 $\frac{l}{4}$ 的长度。B 砖放在水平桌面上，砖的端面与桌边平行。为保持两砖都不翻倒，B 砖伸出桌边的长度 $x$ 的最大值是\xzanswer{C}
\onepicture[w=5cm]{\includesvg{figs/07}}
\fourchoices[answer=C]
{$\frac{l}{8}$}
{$\frac{l}{4}$}
{$\frac{3l}{8}$}
{$\frac{l}{2}$}

%% answer: C
%% memo:
\memoanswer{
根据力矩平衡分析，当 B 砖外移到某处时，A、B 将一起翻倒，故可将 A、B 作为整体考虑，砖即将翻倒时，桌边左侧的砖与桌分离，桌对砖的支持力集中在桌边 O 轴上，使砖转动的力矩由重力引起。A 的重力对 O 轴的力矩 $M_{1}=G\left(x-\frac{l}{4}\right)$，B 的重力对 O 轴的力矩 $M_{2}=G\left(\frac{l}{2}-x\right)$，由 $M_{1}=M_{2}$ 解出 $x=\frac{3l}{8}$，故 C 正确。
}

\item
%% number: 8
%% paperName: 1991年全国高考第 8 题 2 分
%% typeId: 1
%% type: 单选题
%% chapter: 力物体的平衡
%% point: 力矩平衡
%% method: 无
%% score: 2
%% degree: 650
%% duplicateId: 0
%% body:
如图，一均匀木棒 OA 可绕过 O 点的水平轴自由转动。现有一方向不变的水平力 $F$ 作用于该棒的 A 点，使棒从竖直位置缓慢转到偏角 $\theta < 90^{\circ}$ 的某一位置。设 $M$ 为力 $F$ 对转轴的力矩，则在此过程中\xzanswer{B}
\onepicture[w=4cm]{\includesvg{figs/08}}
\fourchoices[answer=B]
{$M$ 不断变大，$F$ 不断变小}
{$M$ 不断变大，$F$ 不断变大}
{$M$ 不断变小，$F$ 不断变小}
{$M$ 不断变小，$F$ 不断变大}

%% answer: B
%% memo:
\memoanswer{
$\theta$ 增大时，重力臂 $l_{2}$ 增大，由 $M_{2}=Gl_{2}$ 可知 $M_{2}$ 增大；因缓慢转动 $v$ 不变，故 $M=M_{2}$，即 $M$ 也将增大，因 $\theta$ 增大时，$F$ 力臂 $l_{1}$ 减小，由 $M=Fl_{1}$ 得 $F$ 增大，故 B 正确。
}

\item
%% number: 9
%% paperName: 1991年全国高考第 9 题 2 分
%% typeId: 1
%% type: 单选题
%% chapter: 电路
%% point: 分压与分流
%% method: 无
%% score: 2
%% degree: 650
%% duplicateId: 0
%% body:
一伏特计由电流表 G 与电阻 $R$ 串联而成，如图所示。若在使用中发现此伏特计的读数总比准确值稍小一些，采用下列哪种措施可能加以改进？\xzanswer{D}
\onepicture[w=5cm]{figs/09.png}
\fourchoices[answer=D]
{在 $R$ 上串联一比 $R$ 小得多的电阻}
{在 $R$ 上串联一比 $R$ 大得多的电阻}
{在 $R$ 上并联一比 $R$ 小得多的电阻}
{在 $R$ 上并联一比 $R$ 大得多的电阻}

%% answer: D
%% memo:
\memoanswer{
读数总稍小于准确值，意味着 $R$ 分去的电压稍大，即 $R$ 值稍大，应略微减少 $R$ 值，采取并联一个大电阻的方法，故 D 正确。
}

\item
%% number: 10
%% paperName: 1991年全国高考第 10 题 2 分
%% typeId: 1
%% type: 单选题
%% chapter: 电场
%% point: 电势能电势差电场力做功
%% method: 无
%% score: 2
%% degree: 650
%% duplicateId: 0
%% body:
两带电小球，电量分别为 $+q$ 和 $-q$，固定在一长度为 $l$ 的绝缘细杆的两端，置于电场强度为 $E$ 的匀强电场中，杆与场强方向平行，其位置如图所示。若此杆绕过 O 点垂直于杆的轴线转过 $180^{\circ}$，则在此转动过程中电场力做的功为\xzanswer{C}
\onepicture[w=5.5cm]{figs/10.png}
\fourchoices[answer=C]
{零}
{$qEl$}
{$2qEl$}
{$\pi qEl$}

%% answer: C
%% memo:
\memoanswer{
电场力做功与路径无关，$+q$ 沿电场力 $F$ 方向移动了 $l$，$-q$ 沿与 $F$ 等大反向电场力作用下移动了 $l$，所以电场力做功为 $2qEl$，故 C 正确。
}

\item
%% number: 11
%% paperName: 1991年全国高考第 11 题 2 分
%% typeId: 1
%% type: 单选题
%% chapter: 功和能
%% point: 动能定理
%% method: 无
%% score: 2
%% degree: 650
%% duplicateId: 0
%% body:
图中 ABCD 是一条长轨道，其中 AB 段是倾角为 $\theta$ 的斜面，CD 段是水平的。BC 是与 AB 和 CD 都相切的一小段圆弧，其长度可以略去不计。一质量为 $m$ 的小滑块在 A 点从静止状态释放，沿轨道滑下，最后停在 D 点。A 点和 D 点的位置如图所示。现用一沿着轨道方向的力推滑块，使它缓慢地由 D 点推回到 A 点时停下。设滑块与轨道间的摩擦系数为 $\mu$，则推力对滑块做的功等于\xzanswer{B}
\onepicture[w=6cm]{figs/11.png}
\fourchoices[answer=B]
{$mgh$}
{$2mgh$}
{$\mu mg\left(s + \frac{h}{\sin\theta}\right)$}
{$\mu mgs + \mu mgh\,\text{ctg}\,\theta$}

%% answer: B
%% memo:
\memoanswer{
由动能定理得，木块从 A 到 D 克服摩擦力做功 $W_f = mgh$，缓慢推木块从 D 到 A 的过程中，动能增量为零，势能增加了 $mgh$，推力做的功为 $W_F$，则 $W_F - W_f = mgh$，解得 $W_F = 2mgh$，故 B 正确。
}

\item
%% number: 12
%% paperName: 1991年全国高考第 12 题 2 分
%% typeId: 1
%% type: 单选题
%% chapter: 电磁感应
%% point: 楞次定律
%% method: 无
%% score: 2
%% degree: 650
%% duplicateId: 0
%% body:
M 和 N 是绕在一个环形铁心上的两个线圈，绕法和线路如图。现将开关 K 从 a 处断开，然后合向 b 处。在此过程中，通过电阻 $R_{2}$ 的电流方向是\xzanswer{A}
\onepicture[w=4cm]{figs/12.png}
\fourchoices[answer=A]
{先由 c 流向 d，后又由 c 流向 d}
{先由 c 流向 d，后由 d 流向 c}
{先由 d 流向 c，后又由 d 流向 c}
{先由 d 流向 c，后由 c 流向 d}

%% answer: A
%% memo:
\memoanswer{
K 接 a 时，由安培定则知铁心中磁感线为顺时针方向；K 从 a 处断开，原电流减弱，线圈 N 中磁通量减少。由楞次定律知，感生电流的磁场要阻碍原磁通量的减少，因此 N 内感生电流的磁感线方向与原磁感线方向相同，为顺时针方向，再由安培定则判定感生电流由 c 向 d；K 接 b 时，M 中电流反向增加，铁心中磁感线为逆时针方向，N 中磁通量增加，由楞次定律知，感生电流的磁场要阻碍原磁通量的增加，因此 N 内感生电流的磁感线方向与原磁感线方向相反，为顺时针方向，再由安培定则判定感生电流由 c 向 d，故 A 正确。
}

\item
%% number: 13
%% paperName: 1991年全国高考第 13 题 2 分
%% typeId: 1
%% type: 单选题
%% chapter: 气体性质
%% point: 液柱移动定性分析
%% method: 无
%% score: 2
%% degree: 650
%% duplicateId: 0
%% body:
两端封闭的等臂 U 形管中，两边的空气柱 a 和 b 被水银柱隔开。当 U 形管竖直放置时，两空气柱的长度差为 $h$，如图所示。现将这个管平放，使两臂位于同一水平面上，稳定后两空气柱的长度差为 $l$，若温度不变则\xzanswer{A}
\onepicture[w=3.5cm]{figs/13.png}
\fourchoices[answer=A]
{$l > h$}
{$l = h$}
{$l = 0$}
{$l < h$，$l \neq 0$}

%% answer: A
%% memo:
\memoanswer{
竖直放置时，$p_a = p_b + \rho gh$，即 $p_a > p_b$；水平放置后，因水银柱在同一水平面上而压强相等；两空气柱压强将趋于平衡，即 $p_a$ 减小，$p_b$ 增大，由玻意耳定律 $pV =$ 恒量，可知 $p_a$ 减小，$V_a$ 会增大，即 $l_a$ 增大，$l_b$ 相应减小，空气柱长度差增加，即 $l > h$，故 A 正确。
}

\item
%% number: 14
%% paperName: 1991年全国高考第 14 题 3 分
%% typeId: 2
%% type: 多选题
%% chapter: 功和能
%% point: 功
%% method: 无
%% score: 3
%% degree: 850
%% duplicateId: 0
%% body:
下列哪些是能量的单位？\xzanswer{ACD}
\fourchoices[answer=ACD]
{焦耳}
{瓦特}
{千瓦小时}
{电子伏特}

%% answer: ACD
%% memo:
\memoanswer{
由题意知，除 B 项的瓦特是功率单位外，其余都是能量单位，故 ACD 正确，B 错误。
}

\item
%% number: 15
%% paperName: 1991年全国高考第 15 题 3 分
%% typeId: 2
%% type: 多选题
%% chapter: 固体液体的性质
%% point: 固体的基本性质
%% method: 无
%% score: 3
%% degree: 850
%% duplicateId: 0
%% body:
下列固态物质哪些是晶体？\xzanswer{ABD}
\fourchoices[answer=ABD]
{雪花}
{黄金}
{玻璃}
{食盐}

%% answer: ABD
%% memo:
\memoanswer{
晶体有一定的熔点，除玻璃不是晶体外，其余都是晶体，故 ABD 正确，C 错误。
}

\item
%% number: 16
%% paperName: 1991年全国高考第 16 题 3 分
%% typeId: 2
%% type: 多选题
%% chapter: 光学
%% point: 光谱
%% method: 无
%% score: 3
%% degree: 650
%% duplicateId: 0
%% body:
关于光谱，下面说法中正确的是\xzanswer{AC}
\fourchoices[answer=AC]
{炽热的液体发射连续光谱}
{太阳光谱中的暗线说明太阳上缺少与这些暗线相应的元素}
{明线光谱和暗线光谱都可用于对物质成分进行分析}
{发射光谱一定是连续光谱}

%% answer: AC
%% memo:
\memoanswer{
炽热的液体发射的光谱为连续光谱，故 A 正确；太阳光谱的暗线是太阳光经地球大气层的吸收而产生的，不是太阳上缺少这些暗线对应的元素，故 B 错误；明线光谱是元素的特征谱线，而暗线光谱由元素吸收产生，故 C 正确；发射光谱包括连续光谱和明线光谱，而不一定是连续光谱，故 D 错误。
}

\item
%% number: 17
%% paperName: 1991年全国高考第 17 题 3 分
%% typeId: 2
%% type: 多选题
%% chapter: 电磁感应
%% point: 电磁感应现象
%% method: 无
%% score: 3
%% degree: 650
%% duplicateId: 0
%% body:
恒定的匀强磁场中有一圆形的闭合导体线圈，线圈平面垂直于磁场方向。当线圈在此磁场中做下列哪种运动时，线圈中能产生感生电流？\xzanswer{CD}
\fourchoices[answer=CD]
{线圈沿自身所在的平面做匀速运动}
{线圈沿自身所在的平面做加速运动}
{线圈绕任意一条直径做匀速转动}
{线圈绕任意一条直径做变速转动}

%% answer: CD
%% memo:
\memoanswer{
A、B 中磁通量没有变化，所以没有感生电流；C、D 中磁通量发生变化，所以有感生电流，故 AB 错误，CD 正确。
}

\item
%% number: 18
%% paperName: 1991年全国高考第 18 题 3 分
%% typeId: 2
%% type: 多选题
%% chapter: 光学
%% point: 光的折射
%% method: 无
%% score: 3
%% degree: 650
%% duplicateId: 0
%% body:
一束光从空气射向折射率 $n = \sqrt{2}$ 的某种玻璃的表面，如图所示。$i$ 代表入射角，则\xzanswer{BCD}
\onepicture[w=5cm]{figs/18.png}
\fourchoices[answer=BCD]
{当 $i > 45^{\circ}$ 时会发生全反射现象}
{无论入射角 $i$ 是多大，折射角 $r$ 都不会超过 $45^{\circ}$}
{欲使折射角 $r = 30^{\circ}$，应以 $i = 45^{\circ}$ 的角度入射}
{当入射角 $i = \arctan\sqrt{2}$ 时，反射光线跟折射光线恰好互相垂直}

%% answer: BCD
%% memo:
\memoanswer{
由发生全反射的条件之一为光线由光密介质射向光疏介质，故 A 错误；由折射定律 $n = \frac{\sin i}{\sin r}$ 知，当折射角 $r$ 为 $45^{\circ}$ 时，入射角 $i$ 为 $90^{\circ}$，所以折射角 $r$ 不会超过 $45^{\circ}$，故 B 正确；当折射角 $r$ 为 $30^{\circ}$ 时，入射角 $i$ 为 $45^{\circ}$，故 C 正确；当 $i = \arctan\sqrt{2}$ 时，$\sin i = \frac{\sqrt{6}}{3}$，$\cos i = \frac{\sqrt{3}}{3}$，由折射定律得 $\sin r = \frac{\sin i}{n} = \frac{\sqrt{3}}{3}$，由三角函数知，$\cos r = \frac{\sqrt{6}}{3}$，故反射光线与折射光线恰好垂直，故 D 正确。
}

\item
%% number: 19
%% paperName: 1991年全国高考第 19 题 3 分
%% typeId: 1
%% type: 单选题
%% chapter: 交流电
%% point: 交流电
%% method: 无
%% score: 3
%% degree: 650
%% duplicateId: 0
%% body:
一矩形线圈，绕垂直于匀强磁场并位于线圈平面内的固定轴转动。线圈中的感生电动势 $e$ 随时间 $t$ 的变化如图所示。下面说法中正确的是\xzanswer{D}
\onepicture[w=6cm]{figs/19.png}
\fourchoices[answer=D]
{$t_{1}$ 时刻通过线圈的磁通量为零}
{$t_{2}$ 时刻通过线圈的磁通量的绝对值最大}
{$t_{3}$ 时刻通过线圈的磁通量变化率的绝对值最大}
{每当 $e$ 变换方向时，通过线圈的磁通量绝对值都为最大}

%% answer: D
%% memo:
\memoanswer{
感生电动势 $e$ 的大小跟穿过线圈的磁通量改变快慢有关，线圈平面垂直于磁场方向，即处于中性面时，磁通量 $\Phi$ 最大，而 $\frac{\Delta\Phi}{\Delta t}=0$，故 $e=0$，即图中 $t_{1}$、$t_{3}$ 时刻，也就是 $e$ 变换方向的时刻，故 ABC 错误，D 正确。
}

\item
%% number: 20
%% paperName: 1991年全国高考第 20 题 3 分
%% typeId: 1
%% type: 单选题
%% chapter: 运动和力
%% point: 变加速直线运动
%% method: 无
%% score: 3
%% degree: 450
%% duplicateId: 0
%% body:
一物体从某一高度自由落下，落在直立于地面的轻弹簧上，如图所示。在 A 点，物体开始与弹簧接触，到 B 点时，物体速度为零，然后被弹回。下列说法中正确的是\xzanswer{C}
\onepicture[w=4cm]{figs/20.png}
\fourchoices[answer=C]
{物体从 A 下降到 B 的过程中，动能不断变小}
{物体从 B 上升到 A 的过程中，动能不断变大}
{物体从 A 下降到 B，以及从 B 上升到 A 的过程中，速率都是先增大，后减小}
{物体在 B 点时，所受合力为零}

%% answer: C
%% memo:
\memoanswer{
在 A 点，物体开始与弹簧接触，弹簧尚未形变，弹力 $N=0$；物体受重力 $G$ 作用继续向下加速，动能增大。物体下落到 O 点时，弹簧形变产生的弹力 $N$ 已逐渐增大到与 $G$ 相等，此时物体所受合力为零，速度达到最大；从 O 到 B 过程中，弹簧形变增大，$N$ 增大，$N>G$，合力 $F=N-G$，与 $v$ 反向，物体做减速运动到 B，动能减小至零，因此物体从 A 到 B 过程中，速率先增大，后减小；同理分析可知，从 B 到 A 亦然，故 C 正确。
}

\item
%% number: 21
%% paperName: 1991年全国高考第 21 题 3 分
%% typeId: 2
%% type: 多选题
%% chapter: 气体性质
%% point: 气体图象
%% method: 无
%% score: 3
%% degree: 650
%% duplicateId: 0
%% body:
一定质量的理想气体经历如图所示的一系列过程，ab、bc、cd 和 da 这四段过程在 $p-T$ 图上都是直线段，其中 ab 的延长线通过坐标原点 O，bc 垂直于 ab，而 cd 平行于 ab。由图可以判断\xzanswer{BCD}
\onepicture[w=5cm]{\includesvg{figs/21}}
\fourchoices[answer=BCD]
{ab 过程中气体体积不断减小}
{bc 过程中气体体积不断减小}
{cd 过程中气体体积不断增大}
{da 过程中气体体积不断增大}

%% answer: BCD
%% memo:
\memoanswer{
由图可知 $\frac{p_{a}}{T_{a}}=\frac{p_{b}}{T_{b}}$，所以 ab 为等容过程，由气态方程 $\frac{pV}{T}=C$（恒量），可得 $\frac{p}{T}=\frac{C}{V}$，对 c、d、a 三点，分别有 $\frac{p_{c}}{T_{c}}=\frac{C}{V_{c}}$，$\frac{p_{d}}{T_{d}}=\frac{C}{V_{d}}$，$\frac{p_{a}}{T_{a}}=\frac{C}{V_{a}}$，由图中各点坐标长度之比可知 $\frac{p_{c}}{T_{c}}>\frac{p_{d}}{T_{d}}>\frac{p_{a}}{T_{a}}$，即 $\frac{C}{V_{c}}>\frac{C}{V_{d}}>\frac{C}{V_{a}}$，所以 $V_{c}<V_{d}<V_{a}$，故 BCD 正确。
}

\item
%% number: 22
%% paperName: 1991年全国高考第 22 题 3 分
%% typeId: 3
%% type: 填空题
%% chapter: 运动和力
%% point: 牛顿定律的应用
%% method: 无
%% score: 3
%% degree: 750
%% duplicateId: 0
%% body:
一物体放在一倾角为 $\theta$ 的斜面上，向下轻轻一推，它刚好能匀速下滑。若给此物体一个沿斜面向上的初速度 $v_{0}$，则它能上滑的最大路程是 \tkanswer[3cm]{$\frac{v_{0}^{2}}{4g\sin\theta}$}。
\onepicture[align=right, w=5cm]{\includesvg{figs/22}}

%% answer: $\frac{v_{0}^{2}}{4g\sin\theta}$
\jdanswer{
$\frac{v_{0}^{2}}{4g\sin\theta}$
}

%% memo:
\memoanswer{
物体匀速下滑表示沿斜面方向的受力平衡，即 $f = G_x = mg\sin\theta$；上滑过程中由动能定理得 $-(f + G_x)s = 0 - \frac{1}{2}mv_0^2$，两式联立解得 $s = \frac{v_0^2}{4g\sin\theta}$。

评分标准：全题 24 分，每小题 3 分。答案正确的，按答案后面括号内的分数给分；答错的，不答的，都给 0 分。
}

\item
%% number: 23
%% paperName: 1991年全国高考第 23 题 3 分
%% typeId: 3
%% type: 填空题
%% chapter: 原子和原子核
%% point: 半衰期
%% method: 无
%% score: 3
%% degree: 750
%% duplicateId: 0
%% body:
两个放射性元素的样品 A 和 B，当 A 有 $\frac{15}{16}$ 的原子核发生了衰变时，B 恰好有 $\frac{63}{64}$ 的原子核发生了衰变。可知 A 和 B 的半衰期之比 $\tau_{A}:\tau_{B} =$ \tkanswer[2cm]{$3:2$}。

%% answer: 3:2
\jdanswer{
$3:2$
}

%% memo:
\memoanswer{
放射性元素每经一个半衰期，就剩下原质量的 $\frac{1}{2}$，经过 $n$ 个半衰期，剩余质量为 $\left(\frac{1}{2}\right)^n$。设 A、B 分别经过 $n_1$、$n_2$ 次半衰期，则 $\left(\frac{1}{2}\right)^{n_1} = 1 - \frac{15}{16}$，$\left(\frac{1}{2}\right)^{n_2} = 1 - \frac{63}{64}$，又 $n_1 = \frac{t}{\tau_A}$，又 $n_2 = \frac{t}{\tau_B}$，由衰变时间相同得 $n_1 \tau_A = n_2 \tau_B$，解得半衰期之比为 $3:2$。
}

\item
%% number: 24
%% paperName: 1991年全国高考第 24 题 3 分
%% typeId: 3
%% type: 填空题
%% chapter: 气体性质
%% point: 阿伏伽德罗常数
%% method: 无
%% score: 3
%% degree: 650
%% duplicateId: 0
%% body:
已知高山上某处的气压为 $0.40\Uatm$，气温为零下 $30\UdC$，则该处每立方厘米大气中的分子数为 \tkanswer[3cm]{$1.2\times10^{19}$}。（阿伏伽德罗常数为 $6.0\times10^{23}\Umoln$，在标准状态下 1 摩尔气体的体积为 $22.4\UL$。）

%% answer: $1.2\times10^{19}$
\jdanswer{
$1.2\times10^{19}$
}

%% memo:
\memoanswer{
由题意可知，1 摩尔气体在标准状态时作为状态 I $(p_{0}, V_{0}, T_{0})$，在题述高山处的状态作为状态 II $(p_{1}, V_{1}, T_{1})$，由理想气体状态方程得

$\frac{p_{0}V_{0}}{T_{0}}=\frac{p_{1}V_{1}}{T_{1}}$，解得 $V_{1}=\frac{p_{0}V_{0}T_{1}}{p_{1}T_{0}}$，

$1\Ucmc$ 的分子数 $n=\frac{N_{A}}{V_{1}}=\frac{N_{A}p_{1}T_{0}}{p_{0}V_{0}T_{1}}$。

其中 $V_{0}=22.4\UL,\ N_{A}=6.0\times10^{23}\Umoln$，

代入数据得 $n=1.2\times10^{19}$。（答 $1\times10^{19}$ 或答数在 $1.0\times10^{19}\sim1.3\times10^{19}$ 范围内的，都给 3 分。）
}

\item
%% number: 25
%% paperName: 1991年全国高考第 25 题 3 分
%% typeId: 4
%% type: 实验题
%% chapter: 光学
%% point: 测定玻璃折射率
%% method: 无
%% score: 3
%% degree: 650
%% duplicateId: 0
%% body:
在测定玻璃的折射率的实验中，对一块两面平行的玻璃砖，用“插针法”找出与入射光线对应的出射光线。现有甲、乙、丙、丁四位同学分别做出如图的四组插针结果。
\begin{enumerate}
	\item
	从图上看，肯定把针插错了的同学是 \tkanswer{乙}。
	\item
	从图上看，测量结果准确度最高的同学是 \tkanswer{丁}。
\end{enumerate}
\fourpicture[numstyle=chinese, align=right, w=2.6cm]{figs/25a.png}{figs/25b.png}{figs/25c.png}{figs/25d.png}

%% answer: 乙，丁
\jdanswer{
\begin{enumerate}
	\item
	乙

	\item
	丁
\end{enumerate}
}

%% memo:
\memoanswer{
本实验中，为了减小误差，入射角应在 $50^{\circ}\sim70^{\circ}$（甲的入射角过小），针距应尽量大些（丙的针距过小），应选用较宽的玻璃砖，还要特别注意界面线与玻璃界面准确重合，故丁同学测量结果准确度最高；乙图中入射光线与射出光线不平行，因此肯定将针插错了。（乙得 1 分，丁得 2 分。）
}

\item
%% number: 26
%% paperName: 1991年全国高考第 26 题 3 分
%% typeId: 3
%% type: 填空题
%% chapter: 电场
%% point: 带电粒子的平衡
%% method: 整体与隔离
%% score: 3
%% degree: 650
%% duplicateId: 0
%% body:
在场强为 $E$、方向竖直向下的匀强电场中，有两个质量均为 $m$ 的带电小球，电量分别为 $+2q$ 和 $-q$。两小球用长为 $l$ 的绝缘细线相连，另用绝缘细线系住带正电的小球悬挂于 O 点而处于平衡状态，如图所示。重力加速度为 $g$。细线对悬点 O 的作用力等于 \tkanswer[2.5cm]{$2mg + qE$}。
\onepicture[align=right, w=3.5cm]{\includesvg{figs/26}}

%% answer: $2mg + qE$
\jdanswer{
$2mg + qE$
}

%% memo:
\memoanswer{
系统平衡时，两个小球可看成整体分析，两个带电球之间的作用力作为内力，整体受到重力、电场力、拉力的作用，由平衡条件得 $F - 2mg - 2qE + qE = 0$，所以 $F = 2mg + qE$。
}

\item
%% number: 27
%% paperName: 1991年全国高考第 27 题 3 分
%% typeId: 3
%% type: 填空题
%% chapter: 电路
%% point: 等效电路
%% method: 无
%% score: 3
%% degree: 650
%% duplicateId: 0
%% body:
如图所示的电路中，三个电阻的阻值相等，电流表 A$_{1}$、A$_{2}$ 和 A$_{3}$ 的内电阻均可忽略，它们的读数分别为 $I_{1}$、$I_{2}$ 和 $I_{3}$，则 $I_{1}:I_{2}:I_{3} =$ \tkanswer[2.5cm]{$3:2:2$}。
\onepicture[align=right, w=5cm]{figs/27.png}

%% answer: 3:2:2
\jdanswer{
$3:2:2$
}

%% memo:
\memoanswer{
画等效电路的方法：等电位法或分电流法，等效电路如图所示。

本题的电路等效于三个电阻并联，每个电阻上电流均为 $I$，总电流 $3I$ 从电源流出，经过表 A$_{1}$，在右端节点分流，经过上面支路 $R$ 的电流等于 $I$，余下 $2I$ 必定经过表 A$_{3}$，同样，在左端，从下面 $R$ 的电流等于 $I$，余下 $2I$ 必经表 A$_{2}$ 流出，汇总到左节点处仍为 $3I$，因此流过三个表的电流之比为 $3:2:2$。

\onepicture[w=6cm]{figs/27_an.jpg}
}

\item
%% number: 28
%% paperName: 1991年全国高考第 28 题 3 分
%% typeId: 3
%% type: 填空题
%% chapter: 磁场
%% point: 带电粒子在磁场中的运动
%% method: 无
%% score: 3
%% degree: 650
%% duplicateId: 0
%% body:
一质量为 $m$、电量为 $q$ 的带电粒子在磁感应强度为 $B$ 的匀强磁场中作圆周运动，其效果相当于一环形电流，则此环形电流的电流强度 $I =$ \tkanswer[3cm]{$\frac{q^{2}B}{2\pi m}$}。

%% answer: $\frac{q^{2}B}{2\pi m}$
\jdanswer{
$\frac{q^{2}B}{2\pi m}$
}

%% memo:
\memoanswer{
洛伦兹力提供向心力得 $qvB = m\frac{v^2}{r}$，解得 $r = \frac{mv}{qB}$，所以周期 $T = \frac{2\pi r}{v} = \frac{2\pi m}{qB}$，等效电流为 $I = \frac{q}{T} = \frac{q^2B}{2\pi m}$。
}

\item
%% number: 29
%% paperName: 1991年全国高考第 29 题 3 分
%% typeId: 7
%% type: 作图题
%% chapter: 机械波
%% point: 波的图象
%% method: 无
%% score: 3
%% degree: 650
%% duplicateId: 0
%% body:
一列简谐波在 $x$ 轴上传播，波速为 50 米/秒。已知 $t = 0$ 时刻的波形图象如图 1 所示，图中 M 处的质点此时正经过平衡位置沿 $y$ 轴的正方向运动。将 $t = 0.5$ 秒时的波形图象画在图 2 上（至少要画出一个波长）。
\twopicture[numstyle=arabic, align=right, w=5cm]{figs/29a.png}{figs/29b.png}

%% answer: 如图
\jdanswer{
如图
}

%% memo:
\memoanswer{
首先根据图中 M 处的质点此时正经过平衡位置沿 $y$ 轴的正方向运动，利用同侧法判断此列波是向左传播，然后从图中读出波长为 $\lambda = 20\Um$，解得周期 $T = \frac{\lambda}{v} = \frac{20}{50}\Us = 0.4\Us$，则 $t = 0.5\Us = 1\frac{1}{4}T$，所以将图 1 中的波形图向左平移 $\frac{1}{4}$ 波长即可。（波形图象至少要画出一个波长，否则不给这 3 分。）

\onepicture[w=6cm]{figs/29_an.jpg}
}

\item
%% number: 30
%% paperName: 1991年全国高考第 30 题 3 分
%% typeId: 4
%% type: 实验题
%% chapter: 电路
%% point: 测电动势与内阻
%% method: 无
%% score: 3
%% degree: 650
%% duplicateId: 0
%% body:
在用电流表和电压表测电池的电动势和内电阻的实验中，所用电流表和电压表的内阻分别为 $0.1\UO$ 和 $1\UkO$。下面分别为实验原理图及所需的器件图。
\onepicture[align=right, w=8cm]{figs/30a.png}
\begin{enumerate}
	\item
	试在下图中画出连线，将器件按原理图连接成实际电路。
	\item
	一位同学记录的 6 组数据见表。试根据这些数据在下图中画出 $U-I$ 图线。根据图线读出电池的电动势 $\varepsilon =$ \tkanswer[2cm]{$1.46\UV$}，根据图线求出电池内阻 $r =$ \tkanswer[2cm]{$0.72\UO$}。
\end{enumerate}

\begin{center}
\begin{tabular}{|c|c|c|c|c|c|c|}
\hline
$I$（A） & 0.12 & 0.20 & 0.31 & 0.32 & 0.50 & 0.57 \\
\hline
$U$（V） & 1.37 & 1.32 & 1.24 & 1.18 & 1.10 & 1.05 \\
\hline
\end{tabular}
\end{center}

\onepicture[align=right, w=6cm]{figs/30b.png}

%% answer: （1）实物图连接如图所示；（2）$U-I$ 图象如图所示；$1.46\pm0.02$，$0.72\pm0.05$
\jdanswer{
\begin{enumerate}
	\item
	实物图连接如图所示。

	\item
	$U-I$ 图象如图所示；$1.46\pm0.02$，$0.72\pm0.05$。
\end{enumerate}
}

%% memo:
\memoanswer{
（1）连接实物图时，电流表、电压表的正、负接线柱要和电源的正、负极一致，即电流从电流表、电压表的正进、负出；连线不要在“空中”交叉，即要在接线柱上连接。

\onepicture[w=7cm]{figs/30_an1.jpg}

（2）$U-I$ 图象如图所示。由闭合电路的欧姆定律知，$E=U+Ir$，即 $U=E-Ir$，故图象的纵截距为电源电动势，图象斜率的绝对值为电源内阻，从图中读出 $E=1.46\UV$，$r=0.72\UO$。

\onepicture[w=5cm]{figs/30_an2.jpg}

评分标准：本题 3 分，接线出现任何错误都不给这 3 分。$\varepsilon=1.46\UV$，$r=0.72\UO$。

评分标准：全题 5 分。正确画得 $U-I$ 图线给 2 分。$U-I$ 图上由各组数据标出的六个点的位置要准确，连直线时第四组数据（$0.32\UA$，$1.18\UV$）标出的点应该舍去不顾。$\varepsilon$ 的答数在 $1.46\pm0.02\UV$ 范围内的都给 1 分。$r$ 的答数在 $0.72\pm0.05\UO$ 范围内的都给 2 分。
}

\item
%% number: 31
%% paperName: 1991年全国高考第 31 题 5 分
%% typeId: 6
%% type: 计算题
%% chapter: 电路
%% point: 电容
%% method: 无
%% score: 5
%% degree: 750
%% duplicateId: 0
%% body:
图中 $\varepsilon = 10\UV$，$R_{1} = 4\UO$，$R_{2} = 6\UO$，$C = 30\UuF$，电池内阻可忽略。
\begin{enumerate}
	\item
	闭合开关 K，求稳定后通过 $R_{1}$ 的电流。
	\item
	然后将开关 K 断开，求这以后流过 $R_{1}$ 的总电量。
\end{enumerate}
\onepicture[align=right, w=4.5cm]{figs/31.png}

%% answer: （1）$I = 1.0\UA$；（2）$\Delta q = 1.2\times10^{-4}\UC$
\jdanswer{
\begin{enumerate}
	\item
	$I = 1.0\UA$

	\item
	$\Delta q = 1.2\times10^{-4}\UC$
\end{enumerate}
}

%% memo:
\memoanswer{
（1）$I=\frac{\varepsilon}{R_{1}+R_{2}}=1.0\UA$ ①

（2）断开前，电容器上电压为 $IR_{2}$，储存的电量为 $q_{1}=CIR_{2}$ ②

断开，待稳定后，电容器上电压为 $\varepsilon$，储存的电量为 $q_{2}=C\varepsilon$ ③

流过 $R_{1}$ 的总电量为

$\Delta q = C(\varepsilon - IR_{2})$ ④

$=1.2\times10^{-4}\UC$

评分标准：本题 5 分。得出①、②、③、④式，各给 1 分。算出数值再给 1 分。
}

\item
%% number: 32
%% paperName: 1991年全国高考第 32 题 5 分
%% typeId: 6
%% type: 计算题
%% chapter: 光学
%% point: 透镜
%% method: 无
%% score: 5
%% degree: 650
%% duplicateId: 0
%% body:
用焦距 $8\Ucm$ 的凸透镜，使一根每小格为 $1\Umm$ 的直尺成像在直径是 $6.4\Ucm$ 的圆形光屏上。要求光屏上显示 16 个小格，应将直尺放在离透镜多远的地方？已知直尺和光屏都垂直于透镜的主光轴，光屏的圆心在主光轴上，直尺与主光轴相交。

%% answer: $10\Ucm$
\jdanswer{
$10\Ucm$
}

%% memo:
\memoanswer{
按题目的要求，在屏上能成像的一段物高 $y = 1.6\Ucm$。屏直径即像高 $y' = 6.4\Ucm$。由放大率公式得

$\frac{v}{u} = \frac{y'}{y}$，解得 $v = 4u$ ①

由成像公式得

$\frac{1}{u} + \frac{1}{v} = \frac{1}{f}$，得 $\frac{1}{u} + \frac{1}{4u} = \frac{1}{f}$ ②

解得 $u = \frac{5}{4}f = \frac{5}{4}\times8\Ucm = 10\Ucm$，

所以直尺到透镜的距离应是 10 厘米。

评分标准：全题 5 分。得出①式给 3 分。得出②式给 1 分。明确表示出直尺到透镜的距离为 10 厘米再给 1 分。
}

\item
%% number: 33
%% paperName: 1991年全国高考第 33 题 8 分
%% typeId: 6
%% type: 计算题
%% chapter: 运动和力
%% point: 动量和能量
%% method: 无
%% score: 8
%% degree: 450
%% duplicateId: 0
%% body:
在光滑的水平轨道上有两个半径都是 $r$ 的小球 A 和 B，质量分别为 $m$ 和 $2m$，当两球心间的距离大于 $l$（$l$ 比 $2r$ 大得多）时，两球之间无相互作用力；当两球心间的距离等于或小于 $l$ 时，两球间存在相互作用的恒定斥力 $F$。设 A 球从远离 B 球处以速度 $v_{0}$ 沿两球连心线向原来静止的 B 球运动，如图所示。欲使两球不发生接触，$v_{0}$ 必须满足什么条件？
\onepicture[align=right, w=6cm]{figs/33.png}

%% answer: $v_{0} < \sqrt{\frac{3F(l-2r)}{m}}$
\jdanswer{
$v_{0} < \sqrt{\frac{3F(l-2r)}{m}}$
}

%% memo:
\memoanswer{
解一：A 球向 B 球接近至 A、B 间的距离小于 $l$ 之后，A 球的速度逐步减小，B 球从静止开始加速运动，两球间的距离逐步减小。当 A、B 的速度相等时，两球间的距离最小。若此距离大于 $2r$，则两球就不会接触。所以不接触的条件是

$v_{1}=v_{2}$ ①

$l+s_{2}-s_{1}>2r$ ②

其中 $v_{1}$、$v_{2}$ 为当两球间距离最小时 A、B 两球的速度；$s_{1}$、$s_{2}$ 为两球间距离从 $l$ 变至最小的过程中，A、B 两球通过的位移。

由牛顿定律得 A 球在减速运动而 B 球作加速运动的过程中，A、B 两球的加速度大小为

$a_{1}=\frac{F}{m},\quad a_{2}=\frac{F}{2m}$ ③

设 $v_{0}$ 为 A 球的初速度，则由匀加速运动公式得

$v_{1}=v_{0}-\frac{F}{m}t,\quad v_{2}=\frac{F}{2m}t,$ ④

$s_{1}=v_{0}t-\frac{1}{2}\cdot\frac{F}{m}t^{2},\quad s_{2}=\frac{1}{2}\cdot\frac{F}{2m}t^{2}$ ⑤

联立解得

$v_{0}<\sqrt{\frac{3F(l-2r)}{m}}$ ⑥

解二：A 球向 B 球接近至 A、B 间的距离小于 $l$ 之后，A 球的速度逐步减小，B 球从静止开始加速运动，两球间的距离逐步减小。当 A、B 的速度相等时，两球间的距离最小。若此距离大于 $2r$，则两球就不会接触。所以不接触的条件是

$v_{1}=v_{2}$ ①

$l+s_{2}-s_{1}>2r$ ②

其中 $v_{1}$、$v_{2}$ 为当两球间距离最小时 A、B 两球的速度；$s_{1}$、$s_{2}$ 为两球间距离从 $l$ 变至最小的过程中，A、B 两球通过的位移。

设 $v_{0}$ 为 A 球的初速度，则由动量守恒定律得

$mv_{0}=mv_{1}+2mv_{2}$ ③

由动能定理得

$Fs_{1}=\frac{1}{2}mv_{0}^{2}-\frac{1}{2}mv_{1}^{2}$ ④

$Fs_{2}=\frac{1}{2}(2m)v_{2}^{2}$ ⑤

联立解得

$v_{0}<\sqrt{\frac{3F(l-2r)}{m}}$ ⑥

评分标准：全题共 8 分。得出①式给 1 分。得出②式给 2 分。若②式中“$>$”写成“$\geq$”的也给这 2 分。在写出①、②两式的条件下，能写出③、④、⑤式，每式各得 1 分。如只写出③、④、⑤式，不给这 3 分。得出结果⑥再给 2 分。若⑥式中“$<$”写成“$\leq$”的也给这 2 分。
}

\end{enumerate}

\end{document}
